% Standalone fragment; input from the main research report.
\subsection{D：完整两轮轨迹的边际校准与最终门控}
\label{sec:D-audited}
一次统计单位是一条独立外生路径及其完整两轮影子轨迹。初始历史有12点，
首轮执行2格后得到14点，第二轮执行剩余6格后得到20点。
四个候选在每轮均参与选择与联合拟合；门控不删候选，前轮被拒候选允许下一轮重新进入。
因此，最终唯一保留与任一轮误拒真候选是两个不同事件。

以取得的局部AC目标 $\widehat L_{i,r,g}$ 定义
\[
 S_i=\max_{r=1,2}\left\{\widehat L_{i,r,h_i}-\min_g\widehat L_{i,r,g}\right\},
 \qquad A_i=\max_{r=1,2}\widehat L_{i,r,h_i}.
\]
任何必要候选拟合失败或漏真均以 $+\infty$ 保留在校准分母中。
64条独立校准轨迹全部完整；两项规则各分配0.05风险预算，
第 $\lceil65\times0.95\rceil=62$ 次序统计量给出
\[
 q_S=4.362284013991641,\qquad q_A=65.99035043496845.
\]
严格超过门限才拒绝。若完整数值策略事前固定，校准与新轨迹可交换，
每项边际秩界为 $3/65$，两项并集界为 $6/65=0.0923077$。
概率同时对随机校准集和新轨迹取平均；该结论不保证给定已实现校准集、
历史、拓扑或任意连续阻抗的条件风险，也不是条件有效的逐轮alpha-spending。
局部拟合损失经过校准仍不是全局MLE或全非线性解释域的分离证书。

96条评价轨迹全部保留在分母中，最终算法唯一正确保留为
$37/96=38.54\%$，95\%两侧Clopper--Pearson区间为
$[28.78\%,49.03\%]$。任一轮真候选误拒为
$10/96=10.42\%$，相应区间为 $[5.11\%,18.32\%]$；
该有限样本没有经验确认实际风险不超过0.1。
最终未决53条、错误单候选2条、空集4条；选择失败、必要联合拟合失败、
实际电压失败及漏真均为 $0/96$，零事件的95\%两侧上界为3.77\%。
最终排除错误候选总数为 $132/288$，候选分母不作为独立统计样本。
首轮误拒3条，第二轮误拒9条，两轮均误拒2条；相对/绝对规则的整轨迹误拒
分别为5/96和6/96，二者并集为10/96。
star、strong cherry和weak cherry的完整分母分别为21、40、35，
唯一正确分别为0、36、1；弱边和star的未决现象仍限制总体恢复。

全轨迹平均耗时26.45726秒，其中两轮选择共25.04336秒，平均研究AC调用
880.7917次；两轮联合拟合平均219.6146次调用。
原经验分析逐方法行只计第二轮选择11.11061秒、382.03125次调用，
不能用该行代表全两轮选择成本。拟合与选择耗时分开归档。
768份保存fit的既有AC目标回放均为verified，最大loss差为
$9.095\times10^{-13}$；紧凑离线核验器复算事件和计时，未重新调用AC或优化器。
H、C和D具有不同历史、数值策略、方法集合及门限，不能将其成功率差解释为
反馈重规划的因果收益。

\subsection{C的总体外压力测试与漏真}
\label{sec:stress-audited}
以下压力批次使用C的冻结数值策略和128条校准轨迹的门限；1条校准必要拟合失败
以 $+\infty$ 保留，门限为8.61031136和75.77255698。
压力分布偏离校准总体，所以下表只报告固定门限下的描述性事件，
不沿用校准总体的误拒保证。各列均按完整独立历史分母统计。
列方法依次是无EV、恒充、固定pulse、随机动作、旧动作库、固定MLE的MILP和refit MILP。

\begin{table}[htbp]
\centering\small
\begin{tabular}{lrrrrrrrr}
\hline
压力 & $N$ & 无EV & 恒充 & pulse & 随机 & 旧库 & 固定MLE & refit\\
\hline
bias：唯一正确 &16&3&9&9&8&8&7&9\\
bias：误拒真候选&16&1&2&2&2&1&2&2\\
source drift：唯一正确&16&2&6&6&5&6&5&6\\
source drift：误拒真候选&16&0&0&0&0&0&0&0\\
load trend：唯一正确&16&0&3&2&1&3&3&3\\
load trend：误拒真候选&16&0&0&0&0&0&0&0\\
missing truth：唯一正确&32&2&3&3&6&6&4&4\\
missing truth：误拒真候选&32&1&0&0&0&0&1&0\\
\hline
\end{tabular}
\caption{压力批次的完整历史事件分子；误拒和漏真分别记账。}
\end{table}

refit MILP在bias、source drift、load trend、missing truth四批的最终排错总数
分别为 $33/48,18/48,9/48,23/74$；相应最终未决为5、10、13、23条。
最后一批有 $10/32$ 真实拓扑不在候选库，所有七方法的漏真事件均为10/32；
refit的误拒为0/32，但真实候选丢失的合并事件仍为10/32。
缺少真拓扑的历史不能由候选内小残差或低误拒数修复，也不能被删出分母。
load trend的fixed pulse有1/16必要拟合失败，其余方法压力批次必要拟合失败为0；
所有四批七方法的选择回退与实际电压失败均为0。
MILP时限终止、缺bound、QP次数另有完整记账；时限终止不等同于数学不可行，
也不等同于实际选择回退。refit的四批MILP时限终止总次数为35、29、32、66。
refit单次选择平均耗时为13.17390、10.91653、12.27065、8.96545秒；
全方法整条历史平均耗时为27.50830、23.96372、27.50788、18.52506秒。
R/X分量误差按历史、轮次、根干/叶边/内部边分别归档，报告完整分母、
适用分母、收敛数和缺失数，未将两轮或逐边值当独立历史合并。

\subsection{配对补充协议与续接检查点}
\label{sec:paired-D-checkpoint}
旧冻结pilot的两条seed（941101、941102）完整通过重放、共同初值和操作合同，
平均16.87101秒。正式32条校准（940001--940032）和48条独立评价
（950001--950048）当时尚未启动。
原八小时预算已结束；用户于2026-10-13 01:54 UTC授权续接后，
另存continuation\_v2协议，将执行截止延长到03:00 UTC，
保留原freeze/pilot字节，只改变截止时间与包导入兼容入口。
统计策略、seed、共同H14的R/X初值、共同fit seed、拟合预算与失败规则保持不变。

续接版重新冻结并重跑相同两条pilot：两条均完整，但平均74.58163秒，
超过事前固定的45秒扩批门槛；因此停止，正式分母为已尝试0/32与0/48，
未形成配对门限或配对效果评价，未绕过门槛，也没有把续接算作原八小时内完成。
这次重跑不增加新的独立pilot单位，pilot也不进入正式分母。
独立只读比对显示：旧/新科学核心源码hash、数值库版本、线程1和BelowNormal合同、
动作、共同初值、全部shadow/paired观测数组、fit损失、proxy目标及AC调用计数一致，
两次内部重放均exact。旧/新wall time不同；CPU负载、缓存和频率没有同期记录，
证据不足以归因具体机器负载原因。
新的16份final fit保存状态目标回放也均verified；该核验不重新优化，
不证明正式配对恢复率、反馈收益或全局MLE。
冻结协议、源hash、失败门槛、分阶段耗时、检查点及可重复读取命令随证据包保留。
